\begin{textsample}{NEG Dim 4 – persona\_gemini – Score -1.00 – t400\_gemini.txt}  \label{ex:f4_neg_006}
The transport sector is responsible for over a quarter of the EU 's greenhouse gas emissions.
In a stark failure of climate policy, these emissions have risen by 28 % since 1990.
Instead of reversing this trend, governments have used post-COVID recovery funds for investments in polluting carmakers and airlines, a move that prioritizes individualism over the societal good.
This path is unsustainable and requires immediate, radical change.
Governments must stop all investments in new airports and highways and end \textbf{tax} breaks that benefit cars.
Short-haul flights must be banned wherever a viable alternative exists, and sales of new combustion or hybrid cars must be halted by 2028.
The funds currently propping up polluting industries should be redirected to achieve decarbonization by 2040.
This means investing in affordable trains, cycling infrastructure, and public transport.
As we reshape our transport systems, it is essential to support worker transitions into green jobs.
Public action is critical.
Individuals can contribute by contacting their representatives to demand change and by tracking their own personal emissions to understand their impact.

% matched lemmas: tax
\end{textsample}
